\chapter{References}
\label{chap:references}
\begin{thebibliography}{99}
\bibitem{rahman2022unintentional}
M. M. H. Rahman, M. A. Al Naeem, and A. Abubakar,
``Threats From Unintentional Insiders: An Assessment of an
Organization's Readiness Using Machine Learning,''
\emph{IEEE Access}, vol. 10, pp. 110294--110308, 2022.
doi: \url{https://doi.org/10.1109/ACCESS.2022.3214819}.

\bibitem{alshehari2023isolation}
T. Al-Shehari, M. Al-Razgan, T. Alfakih, R. A. Alsowail, and
S. Pandiaraj, ``Insider Threat Detection Model Using Anomaly-Based
Isolation Forest Algorithm,'' \emph{IEEE Access}, vol. 11,
pp. 118170--118185, 2023.
doi: \url{https://doi.org/10.1109/ACCESS.2023.3326750}.

\bibitem{dimitriadis2023fronesis}
A. Dimitriadis, E. Lontzetidis, B. Kulvatunyou, N. Ivezic,
D. Gritzalis, and I. Mavridis, ``Fronesis: Digital Forensics-Based
Early Detection of Ongoing Cyber-Attacks,'' \emph{IEEE Access},
vol. 11, pp. 728--743, 2023.
doi: \url{https://doi.org/10.1109/ACCESS.2022.3233404}.

\bibitem{liu2008isolation}
F. T. Liu, K. M. Ting, and Z.-H. Zhou, ``Isolation Forest,'' in
\emph{Proceedings of the 2008 Eighth IEEE International Conference on
Data Mining}, 2008, pp. 413--422.
doi: \url{https://doi.org/10.1109/ICDM.2008.17}.

\bibitem{cert}
Software Engineering Institute, Carnegie Mellon University,
\emph{CERT Insider Threat Test Dataset, Version 4.2}.
Available: \url{https://www.sei.cmu.edu/library/insider-threat-test-dataset/}.
Accessed: 25 October 2026.

\bibitem{nist53}
Joint Task Force,
\emph{Security and Privacy Controls for Information Systems and Organizations},
NIST Special Publication 800-53, Revision 5, National Institute of
Standards and Technology, September 2020.
doi: \url{https://doi.org/10.6028/NIST.SP.800-53r5}.

\bibitem{nist61}
P. Cichonski, T. Millar, T. Grance, and K. Scarfone,
\emph{Computer Security Incident Handling Guide},
NIST Special Publication 800-61, Revision 2, National Institute of
Standards and Technology, August 2012.
doi: \url{https://doi.org/10.6028/NIST.SP.800-61r2}.

\bibitem{sklearn}
F. Pedregosa et al.,
``Scikit-learn: Machine Learning in Python,''
\emph{Journal of Machine Learning Research}, vol. 12, pp. 2825--2830,
2011. Available: \url{https://jmlr.org/papers/v12/pedregosa11a.html}.

\bibitem{pandas}
The pandas development team,
\emph{pandas documentation}.
Available: \url{https://pandas.pydata.org/docs/}.
Accessed: 25 October 2026.

\bibitem{fastapi}
S. Ramirez,
\emph{FastAPI documentation}.
Available: \url{https://fastapi.tiangolo.com/}.
Accessed: 25 October 2026.

\bibitem{react}
Meta Open Source,
\emph{React documentation}.
Available: \url{https://react.dev/}.
Accessed: 25 October 2026.

\bibitem{python}
Python Software Foundation,
\emph{Python 3 documentation}.
Available: \url{https://docs.python.org/3/}.
Accessed: 25 October 2026.

\bibitem{nistai}
National Institute of Standards and Technology,
\emph{Artificial Intelligence Risk Management Framework (AI RMF 1.0)},
NIST AI 100-1, January 2023.
doi: \url{https://doi.org/10.6028/NIST.AI.100-1}.
\end{thebibliography}\label{chap:last-reference-page}


\chapter{Appendix}
\label{chap:appendix}

\section{Minimum software requirements}
\label{app:software}
\begin{table}[H]
\centering
\tablefont
\caption{Software needed for the current demonstration}
\label{tab:software-requirements}
\begin{tabularx}{\textwidth}{|p{0.25\textwidth}|p{0.28\textwidth}|Y|}
\hline
\textbf{Software} & \textbf{Purpose} & \textbf{Project note}\\\hline
Python 3 & Model scoring, API, and alternative Streamlit UI. & Use an environment compatible with the installed scientific packages.\\\hline
Python dependencies & pandas, NumPy, scikit-learn, joblib, FastAPI, python-multipart, Streamlit, Uvicorn. & Install from the repository requirements file and record resolved versions.\\\hline
Node.js and npm & React and Vite development client. & Use the frontend package manifest and lock file when available.\\\hline
Modern browser & Select local activity files and render the dashboard. & Directory upload support may differ by browser.\\\hline
LaTeX distribution & Compile this report with pdfLaTeX. & The report uses standard packages including TikZ and newtx.\\\hline
Model artifact & Persisted estimator bundle. & Expected path is \code{sprint1\_output/insider\_investigation\_model.joblib}.\\\hline
\end{tabularx}
\end{table}

\section{Minimum hardware requirements}
\label{app:hardware}
The repository does not contain a measured minimum hardware
benchmark. Actual resource use depends on input size, concurrent
processes, Python package versions, and the fitted estimators. The
following table is a practical starting point for a student
demonstration and must not be treated as a tested guarantee.

\begin{table}[H]
\centering
\tablefont
\caption{Practical workstation guidance, not a measured benchmark}
\label{tab:hardware-requirements}
\begin{tabularx}{\textwidth}{|p{0.25\textwidth}|p{0.25\textwidth}|Y|}
\hline
\textbf{Resource} & \textbf{Suggested starting point} & \textbf{Reason}\\\hline
Processor & Modern multi-core desktop or laptop CPU. & CSV parsing and model scoring benefit from general-purpose compute; exact runtime is unmeasured.\\\hline
Memory & 16 GB RAM recommended for the full synthetic dataset. & Chunked input lowers read size, but merged employee-day tables and models remain in memory.\\\hline
Storage & Several gigabytes of free space for source logs, environments, outputs, and report assets. & Dataset copies, notebook artifacts, and dependencies occupy disk space.\\\hline
Display & 1366 by 768 pixels or higher. & Gives reasonable room for the dashboard and report figures.\\\hline
Network & Not required for local analysis after dependencies are installed. & The documented demo binds locally; package installation may require connectivity.\\\hline
\end{tabularx}
\end{table}

\section{Input contract and data lineage}
\label{app:data-lineage}
\begin{table}[H]
\centering
\tablefont
\caption{Input contract checklist}
\label{tab:input-contract}
\begin{tabularx}{\textwidth}{|p{0.24\textwidth}|p{0.36\textwidth}|Y|}
\hline
\textbf{Check} & \textbf{Expected condition} & \textbf{Operator action}\\\hline
Required files & All five expected filenames are present. & Confirm that the selected folder is the intended dataset.\\\hline
Date and user & Each input contains parsable date and user fields. & Correct the source export; do not silently substitute another identity column.\\\hline
Time coverage & Sources cover a consistent period and time zone. & Review date ranges and document any intentional gaps.\\\hline
Duplicates & Cleaning results are understood. & Consider whether exact global deduplication is required.\\\hline
Optional names & Lookup has a recognised user and name column. & Treat display names as sensitive and keep them out of model features.\\\hline
Model schema & Bundle feature list matches the intended model release. & Verify bundle provenance and avoid untrusted serialized files.\\\hline
\end{tabularx}
\end{table}

\section{Test cases and verification plan}
\label{app:test-plan}
\begin{longtable}{|p{0.10\textwidth}|p{0.24\textwidth}|p{0.32\textwidth}|p{0.27\textwidth}|}
\caption{Functional and data-quality test plan}\label{tab:test-cases}\\
\hline
\textbf{ID} & \textbf{Scenario} & \textbf{Expected result} & \textbf{Execution status}\\\hline
\endfirsthead
\hline
\textbf{ID} & \textbf{Scenario} & \textbf{Expected result} & \textbf{Execution status}\\\hline
\endhead
\hline
\multicolumn{4}{r}{Continued on next page}\\
\endfoot
\hline
\endlastfoot
T-01 & Valid folder with five source files and model bundle. & Analysis returns employee summaries and activity rows. & Not recorded as executed in this report.\\\hline
T-02 & Missing one required CSV. & API reports the missing filename; no successful result is returned. & Not recorded as executed in this report.\\\hline
T-03 & Duplicate accepted filename in uploaded folder. & API returns a clear duplicate-file error. & Not recorded as executed in this report.\\\hline
T-04 & CSV missing date or user column. & Scorer raises a schema error identifying the source. & Not recorded as executed in this report.\\\hline
T-05 & Invalid date and blank user rows. & Invalid rows do not contribute to daily aggregates. & Cleaning summaries show no such rows for saved dataset; synthetic fixture test not recorded.\\\hline
T-06 & Model bundle lacks a required model or feature list. & Bundle validation raises an actionable error. & Not recorded as executed in this report.\\\hline
T-07 & Required model feature absent from aggregate. & Inference inserts a zero column in the saved feature order. & Behaviour present in source; automated test not recorded.\\\hline
T-08 & Request employee detail before an analysis. & API returns conflict status and explanatory detail. & Not recorded as executed in this report.\\\hline
T-09 & Request an unknown employee after analysis. & API returns not-found status. & Not recorded as executed in this report.\\\hline
T-10 & API unavailable when UI loads. & UI communicates local service unavailability and retries health. & Interaction visible in source; browser test not recorded.\\\hline
T-11 & Run analysis with a large CSV. & Chunk processing completes within tested memory and time bounds. & Benchmark not recorded; measure before operational use.\\\hline
T-12 & Repeat analysis with the same verified inputs and bundle. & Outputs are reproducible within expected numeric tolerance. & End-to-end repeatability test not recorded.\\\hline
\end{longtable}

\section{Evaluation metrics and saved counts}
\label{app:metrics}
\begin{table}[H]
\centering
\tablefont
\caption{Saved one-percent test confusion counts by model}
\label{tab:all-confusion}
\begin{tabularx}{\textwidth}{|p{0.27\textwidth}|R|R|R|R|}
\hline
\textbf{Model} & \textbf{TN} & \textbf{FP} & \textbf{FN} & \textbf{TP}\\\hline
Logistic Regression & 65,169 & 660 & 287 & 2\\\hline
Decision Tree & 65,247 & 582 & 209 & 80\\\hline
Random Forest & 65,262 & 567 & 194 & 95\\\hline
Isolation Forest & 65,185 & 644 & 271 & 18\\\hline
\end{tabularx}
\end{table}

The confusion counts are useful for reconstructing capacity metrics.
For Random Forest, \(TP+FP=662\), equal to the selected one-percent
queue size. There are \(TP+FN=289\) positive test labels. This yields
precision \(95/662\), recall \(95/289\), and the corresponding F1.
Small differences in displayed values arise from rounding.

\section{Troubleshooting guide}
\begin{table}[H]
\centering
\tablefont
\caption{Common demonstration issues}
\label{tab:troubleshooting}
\begin{tabularx}{\textwidth}{|p{0.27\textwidth}|p{0.32\textwidth}|Y|}
\hline
\textbf{Symptom} & \textbf{Likely cause} & \textbf{Recommended check}\\\hline
Health reports model missing & Bundle absent or wrong path. & Run the notebook model-bundle creation step and verify the expected filename.\\\hline
Browser cannot reach the API & API is stopped or listening on a different address. & Start the API using the documented local command and inspect its terminal.\\\hline
Analysis says files are missing & Folder selection omitted a source or names differ. & Confirm all five exact filenames are present in the selected folder.\\\hline
Scoring fails on a CSV & Unsupported schema, encoding, or malformed content. & Check headers, delimiter, date values, file integrity, and service error detail.\\\hline
No employee details appear & Analysis did not complete or selected user is absent. & Re-run analysis, then select a user returned in the current report.\\\hline
Results differ from prior run & Input batch, model bundle, package version, or row order differs. & Record data and artifact hashes, environment versions, and run settings.\\\hline
\end{tabularx}
\end{table}

\section{Responsible-use checklist}
\label{app:responsible-use}
Before any review involving real people, an accountable organisation
should verify that all of the following are satisfied:
\begin{itemize}
  \item The review purpose, authority, data sources, and data period are
  documented and approved.
  \item Access to raw logs, derived reports, and employee-name lookups
  is limited to authorised personnel.
  \item The model's limitations and relative score meaning are
  explained to each reviewer.
  \item A score cannot, by itself, trigger disciplinary, employment,
  access, or law-enforcement action.
  \item Reviewers can inspect original records, record alternative
  explanations, and correct erroneous identity or source data.
  \item Retention, deletion, disclosure, and audit requirements are
  defined before data is loaded.
  \item The model and thresholds have been assessed for the intended
  population and operating environment.
  \item The organisation has a process to pause the tool when data
  quality, access control, or model behaviour is unacceptable.
\end{itemize}

\section{Glossary}
\begin{longtable}{|p{0.25\textwidth}|p{0.68\textwidth}|}
\caption{Terms used in the report}\label{tab:glossary}\\
\hline
\textbf{Term} & \textbf{Meaning in this project}\\\hline
\endfirsthead
\hline
\textbf{Term} & \textbf{Meaning in this project}\\\hline
\endhead
\hline
\endlastfoot
Activity event & One row in one of the input CSV source logs.\\\hline
Employee-day & Aggregated feature record for one user identifier on one calendar day.\\\hline
Incident-window label & Positive label when a user-day overlaps a listed synthetic incident interval.\\\hline
Review capacity & Fraction or count of highest-ranked rows selected for human review.\\\hline
Average Precision & Ranking metric summarising precision across recall changes as the score threshold varies.\\\hline
Percentile rank & Relative position of a score within the current scored batch.\\\hline
Model agreement & Count of model employee-level ranks that meet the per-model top-ten-percent rule.\\\hline
False positive & Selected test record without a matching positive label under the benchmark rule.\\\hline
False negative & Positive benchmark record not selected at the chosen queue capacity.\\\hline
Data lineage & Trace from raw source records through cleaning and features to scores and outputs.\\\hline
\end{longtable}

\section{Reproducibility checklist}
\begin{enumerate}
  \item Record the repository commit, Python version, dependency
  versions, and operating system.
  \item Verify the dataset release, input filenames, checksums, and
  licensing conditions.
  \item Re-run the notebook from a clean environment and confirm cell
  execution order.
  \item Confirm the user-group split has no overlap and that each
  partition contains positive examples.
  \item Save the training and validation selection outputs, model
  parameters, feature list, and augmentation decision.
  \item Keep the test metrics separate from all selection decisions.
  \item Record the model-bundle hash and the source logs used for each
  dashboard run.
  \item Compile this document twice and inspect page count, contents,
  tables, figures, code listings, and page-number transitions.
\end{enumerate}

\section{Notebook source listing}
\label{app:notebook-source}
The following appendix listing is generated from the project's
notebook code cells. It is included as a reproducibility aid and should
be refreshed from \code{notebook/forensic.ipynb} when the notebook
changes. It is source documentation, not a substitute for executing
the notebook and reviewing its outputs.

\clearpage
\lstinputlisting[caption={Notebook code cells for data preparation, model evaluation, and bundle creation},label={lst:notebook}]{notebook_code_listing.txt}

\section{Scoring core source listing}
\label{app:core-source}
This listing captures the Python functions used by the API and
Streamlit interface to validate the model bundle, read activity data,
create employee-day features, calculate ranks, and format activity
details.

\clearpage
\lstinputlisting[caption={Forensic core implementation},label={lst:core}]{forensic_core_code_listing.txt}

\section{FastAPI source listing}
\label{app:api-source}
This listing documents the local HTTP interface, file staging,
progress state, API errors, and employee detail response.

\clearpage
\lstinputlisting[caption={FastAPI service implementation},label={lst:api}]{api_code_listing.txt}

\section{Streamlit source listing}
\label{app:streamlit-source}
The Streamlit interface is retained as an alternative local
demonstration. Its output is also review-oriented and shows the model
agreement and aggregate activity for selected users.

\clearpage
\lstinputlisting[caption={Alternative Streamlit interface},label={lst:streamlit}]{app_code_listing.txt}

\section{React application source listing}
\label{app:react-source}
The React client implements service health checks, directory upload,
progress polling, report rendering, employee selection, search, and
filtering. The listing is included to show how the browser workflow
relates to the API contract described in Chapter 4.

\clearpage
\lstinputlisting[caption={React dashboard application},label={lst:react}]{react_code_listing.txt}

\section{Interface style source listing}
\label{app:style-source}
The frontend styling is maintained separately from the React
component. These listings record the visual layout and the current
theme overrides used by the saved UI screenshots.

\clearpage
\lstinputlisting[caption={Dashboard base styles},label={lst:styles}]{styles_code_listing.txt}

\clearpage
\lstinputlisting[caption={Dashboard theme styles},label={lst:premium-styles}]{theme_code_listing.txt}
\label{app:end-content}
